\chapter*{Introduction}

% Ajouter manuellement cette section à la table des matières
\addcontentsline{toc}{chapter}{Introduction}

Dans un contexte réglementaire exigeant et en constante évolution, les organismes d'assurance doivent être en mesure de valoriser leurs engagements de manière fiable, tout en assurant la compréhension et la traçabilité des résultats produits. Depuis l'entrée en vigueur de la directive Solvabilité~II, le calcul des provisions techniques repose sur une approche économique tenant compte des conditions de marché. Le Best Estimate occupe une place centrale dans cette valorisation puisqu'il correspond à l'estimation des flux futurs liés aux engagements d'assurance.

En assurance vie, le calcul du Best Estimate peut être réalisé à partir d'une approche déterministe ou d'une approche stochastique. L'approche déterministe repose sur une projection fondée sur un scénario économique central. Elle présente l'avantage d'être plus rapide à mettre en œuvre et facilite l'analyse des résultats. L'approche stochastique repose, quant à elle, sur un ensemble de scénarios économiques permettant de prendre en compte l'incertitude liée aux évolutions des marchés financiers.

Cette seconde approche permet notamment de mieux intégrer les effets liés aux options et garanties présentes dans les contrats d'assurance vie. Elle est cependant plus complexe à mettre en œuvre et nécessite des temps de calcul plus importants.

Les résultats obtenus à partir des approches déterministe et stochastique peuvent donc présenter des différences. Ces écarts dépendent à la fois de l'environnement économique, des caractéristiques du portefeuille et des hypothèses utilisées dans le modèle. Leur compréhension constitue un enjeu important, car elle permet de mieux analyser les résultats produits et de renforcer leur justification dans un cadre réglementaire.

L'objectif de ce mémoire est ainsi de rationaliser le passage du Best Estimate déterministe au Best Estimate stochastique. Il ne s'agit pas de remplacer le calcul stochastique, mais de mieux comprendre les facteurs expliquant les différences observées entre les deux méthodes de valorisation. Une attention particulière est portée à l'identification des variables ayant le plus d'influence sur ces écarts.

Pour répondre à cette problématique, une base de données est construite à partir des résultats d'un modèle ALM interne. Différentes hypothèses économiques et contractuelles sont modifiées de manière contrôlée afin de créer plusieurs configurations de portefeuille. Pour chacune de ces configurations, un Best Estimate déterministe et un Best Estimate stochastique sont calculés.

Dans le cadre de ce mémoire, nous aurons recours à plusieurs modèles de machine learning afin d'évaluer leur capacité à reproduire et à expliquer l'écart entre le Best Estimate déterministe et le Best Estimate stochastique. L'objectif n'est pas uniquement d'obtenir de bonnes performances prédictives. Il est également nécessaire de comprendre le fonctionnement des modèles et d'identifier les variables les plus significatives.

Le mémoire est organisé en plusieurs parties. La première partie présente le cadre réglementaire de Solvabilité~II et les principes de calcul du Best Estimate. La deuxième partie décrit le modèle ALM interne utilisé ainsi que la méthodologie de construction de la base de données. La troisième partie présente les modèles de machine learning retenus et la méthodologie d'apprentissage. La quatrième partie est consacrée à l'analyse des résultats et à l'identification des principaux facteurs expliquant les écarts observés. Enfin, la dernière partie présente les limites de l'étude ainsi que les perspectives d'amélioration possibles.